\section{Por qué estudiar el aprendizaje por refuerzo profundo}

\subsection{Más allá del aprendizaje supervisado}

Muchas aplicaciones de IA del mundo real no son simples mapeos de entrada a salida. El expositor da varios escenarios clave:

\begin{enumerate}
    \item \textbf{Las decisiones tienen consecuencias}: por ejemplo, después de que Spotify recomienda una canción, la siguiente recomendación debe evitar repetir la misma canción y debe mantener continuidad de estilo.
    \item \textbf{No se dispone de supervisión directa}: por ejemplo, en los asistentes de codificación (Cursor, Copilot), el usuario solo puede dar retroalimentación de "bueno" o "malo", en lugar de proporcionar directamente la salida correcta.
    \item \textbf{El objetivo no es diferenciable}: muchos objetivos reales (como la satisfacción del usuario) no se pueden optimizar directamente mediante el gradiente.
\end{enumerate}

\begin{figure}[H]
\centering
\includegraphics[width=0.9\textwidth]{slides-images/slide-18.jpg}
\caption{El curso usa la perspectiva de que "el despliegue cambia las observaciones futuras" para explicar por qué los problemas de decisión secuencial no pueden reducirse a un aprendizaje supervisado ordinario (Slide 18).}
\end{figure}

\subsection{Aplicación amplia en sistemas de IA de alto desempeño}

El expositor presentó una gran cantidad de aplicaciones reales del RL profundo:

\begin{itemize}
    \item \textbf{Robótica}: la marcha del robot cuadrúpedo de Unitree, la manipulación con $\pi_0$ de Physical Intelligence, Google Gemini Robotics.
    \item \textbf{Juegos}: la jugada 37 de AlphaGo mostró la capacidad del RL de descubrir estrategias nuevas que los humanos nunca habían imaginado.
    \item \textbf{Modelos de lenguaje grandes}: casi todos los LLM modernos usan alguna forma de RL para el post-training, sobre todo en el razonamiento avanzado.
    \item \textbf{Otros}: control de tráfico, seguimiento de prompt en modelos generativos de imágenes, diseño de chips (Google TPU).
\end{itemize}

\begin{figure}[H]
\centering
\includegraphics[width=0.88\textwidth]{slides-images/slide-19.jpg}
\caption{El RL profundo ya se ha convertido en un paradigma de entrenamiento importante para sistemas robóticos complejos, especialmente adecuado para tareas físicas que requieren control en lazo cerrado y adaptación en línea (Slide 19).}
\end{figure}

\begin{importantbox}{La capacidad de descubrimiento del RL}
El aprendizaje por refuerzo no solo imita el comportamiento ya presente en los datos, sino que también puede descubrir soluciones completamente nuevas mediante el "ensayo y error". La jugada 37 de AlphaGo es un caso clásico: fue una jugada que ningún jugador humano de go había hecho antes.
\end{importantbox}

\begin{figure}[H]
\centering
\includegraphics[width=0.88\textwidth]{slides-images/slide-22.jpg}
\caption{El curso también incluye explícitamente el post-training de LLM dentro del mapa de aplicaciones del RL profundo, lo que muestra que el RL se ha extendido de la robótica y los juegos a los sistemas de modelos de lenguaje (Slide 22).}
\end{figure}

\subsection{La naturaleza del aprendizaje}

Desde una perspectiva filosófica, aprender de la experiencia parece ser un componente fundamental de la inteligencia. Contar con algoritmos capaces de practicar por su cuenta y mejorar es esencial para construir sistemas verdaderamente inteligentes.

\subsection{Resumen de la sección}

Las razones para estudiar el RL profundo incluyen: (1) muchos problemas están fuera del alcance del aprendizaje supervisado; (2) los sistemas de IA de alto desempeño usan ampliamente el RL; (3) aprender de la experiencia es una capacidad fundamental de la inteligencia; (4) todavía hay una gran cantidad de problemas de investigación abiertos en el área.

